% trace_spaces.tex
%
% Purpose:
%   Define Rayleigh wall coordinates and distinguish the ambient trace space
%   from physically admissible incoming/outgoing trace subspaces.
%
% Main algorithm:
%   Introduces the truncated Rayleigh basis and the Cauchy coordinate vectors
%   D^\pm and N^\pm used throughout the matching system.
%
% Based on:
%   notes/theory/scat_formulation.md and notes/theory/mode_selection.md.
%
% Main changes:
%   The explanatory material is rewritten in English with a compact notation.
%
% Numerical goal:
%   Make explicit what the truncation parameter M controls and what it does not
%   select physically.

\section{Rayleigh Trace Coordinate Spaces}

On any vertical port \(x=X\), the transverse quasiperiodicity suggests the
Rayleigh basis
\[
  \beta_m=\beta+\frac{2\pi m}{d},
  \qquad
  \psi_m(y)=\frac{1}{\sqrt d}\ee^{\ii\beta_m y},
  \qquad
  m\in\mathbb Z .
\]
The longitudinal Rayleigh wavenumber is
\[
  \gamma_m=\sqrt{k^2-\beta_m^2},
  \qquad
  \operatorname{Im}\gamma_m\ge 0 .
\]
The truncated index set is \(|m|\le M\), and
\[
  K=2M+1 .
\]

For each port \(\Gamma_\pm\), write the Dirichlet and outward-normal derivative
trace coefficients as
\[
  D^\pm\in\CC^K,
  \qquad
  N^\pm\in\CC^K .
\]
At the continuous level, these are coordinates for Cauchy traces in
\[
  H^{1/2}(\Gamma_\pm)\times H^{-1/2}(\Gamma_\pm).
\]
After truncation, the ambient coordinate space is
\[
  \Tspace_M=\CC_D^K\times\CC_N^K .
\]

This ambient space is only a coordinate space.  It does not say that an
arbitrary vector \((D,N)\in\Tspace_M\) is the trace of a physical half-lead
solution.  Physical incoming and outgoing data lie in subspaces,
\[
  \Ein^\pm,\Eout^\pm
  \subset
  \CC^K\times\CC^K,
\]
selected by the periodic half-lead radiation or admissibility condition.

\subsection{Rayleigh coordinates versus Bloch modes}

The Rayleigh basis is a wall trace coordinate basis.  A Bloch mode is a one-cell
periodic-lead propagation eigenmode, and in a nonempty periodic cell it is not
usually a single Rayleigh order.  Its wall trace can still be expanded in the
Rayleigh basis, which is why the Bloch-selected subspaces can be represented by
matrices
\[
  \begin{bmatrix}
  \Din^\pm\\
  \Nin^\pm
  \end{bmatrix},
  \qquad
  \begin{bmatrix}
  \Dout^\pm\\
  \Nout^\pm
  \end{bmatrix}.
\]
Local right-going and left-going Rayleigh amplitudes in a center cell are also
different objects from physical incoming and outgoing half-lead modes.  The
labels ``right-going'' and ``left-going'' describe the elementary factors
\(\ee^{\pm \ii\gamma_m x}\) in a homogeneous center region.  Physical incoming
and outgoing behavior in a periodic half-lead is selected by the Bloch
multiplier, and, for unit-circle modes, by an energy-flux or limiting absorption
criterion.
